\section{Assumption sensitivity and identification limits}
\label{sec:sensitivity}

The reconstruction in Section~\ref{sec:comparison} rests on choices that no filing
constrains. This section puts a number on each of them, and the result is not
symmetric: the asset-side choices move the answer by a few points, and the deposit
choice moves it by more than the entire exposure being measured.

Before any of it, the nature of the uncertainty. Everything below is a deterministic
re-solve of the model with one input changed. These are not sampling distributions
and they do not become confidence intervals when read quickly. Where a range appears
it is the range across assumptions a reviewer might reasonably argue for, and the
only standard errors anywhere in this paper are the Monte Carlo errors of the pricing
engine in Appendix~\ref{app:pricing}.

\subsection{The asset side}

Thirteen variants re-run the out-of-sample test of Section~\ref{sec:comparison} with
a single parameter changed and the random seed held fixed, so that differences
between rows measure the assumption rather than simulation noise.
Table~\ref{tab:sweep} reports them.

\begin{table}[htbp]
\centering
\caption{Predicted SVB held-to-maturity price under each assumption variant, against
the 83.40 implied by disclosure. Run at 4,000 paths with common random numbers, so
the base case here is 80.35 rather than the 80.36 of
Table~\ref{tab:oos}; the two are not interchangeable and the difference is
simulation noise, not a result.}
\label{tab:sweep}
\small
\begin{tabular}{lrr}
\toprule
Variant & Predicted & Error (\%) \\
\midrule
Base case & 80.35 & $-3.66$ \\
\addlinespace
Over-10y bullet maturity 15y & 81.43 & $-2.36$ \\
Over-10y bullet maturity 25y & 79.52 & $-4.65$ \\
Pass-through life 20y & 81.48 & $-2.30$ \\
Pass-through life 28y & 79.80 & $-4.33$ \\
Prepayment floor 2\% / ceiling 60\% & 78.30 & $-6.11$ \\
Prepayment floor 8\% / ceiling 35\% & 81.90 & $-1.80$ \\
Prepayment curve twice as steep & 80.71 & $-3.22$ \\
Servicing and guarantee fee 25bp & 80.51 & $-3.47$ \\
Servicing and guarantee fee 75bp & 80.14 & $-3.91$ \\
\addlinespace
Mean reversion half-life 14y & 80.63 & $-3.32$ \\
Mean reversion half-life 4.6y & 80.12 & $-3.94$ \\
Short-rate volatility 100bp & 81.04 & $-2.83$ \\
\bottomrule
\end{tabular}
\end{table}

Three things follow.

The prediction spans 3.60 price points across every variant, from 78.30 to 81.90, and
the error never leaves the interval $[-6.11\%, -1.80\%]$. So the headline result is
not an artefact of one fortunate parameter choice, but neither is it precise to better
than a couple of points.

Every variant errs in the same direction. That is the more informative fact. A model
whose remaining error were noise would straddle the disclosed figure under at least
some plausible assumption, and none of these thirteen does. The bias is structural.

And it does not live in the rate model. The three variants that move the short-rate
specification -- mean reversion half-lives of 14 and 4.6 years against the 7.8 used,
and volatility raised from 62 to 100 basis points -- span 0.92 price points between
them. The variants that change the shape of the asset book span 3.60. Whatever is
left over after calibration is a feature of how contractual maturity buckets and
prepayment behaviour are approximated, not of how the short rate is modelled. This
matters because the short-rate calibration is the part of the model with the weakest
statistical foundation, as the next subsection shows, and it turns out to be the part
that matters least.

\subsection{Identification limits in the short-rate calibration}

Mean reversion and volatility are estimated from the realised path of the three-month
bill rather than from a swaption surface. Table~\ref{tab:ou} runs that estimator
across four windows and three sampling frequencies.

\begin{table}[htbp]
\centering
\caption{Ornstein--Uhlenbeck estimates of the three-month bill across windows and
sampling steps. The specification used throughout this paper is monthly data from
1990, marked. Half-lives are in years, volatility in basis points, long-run mean in
per cent.}
\label{tab:ou}
\small
\begin{tabular}{llrrrrr}
\toprule
Window & Step & Obs. & Half-life & $\sigma$ & Long-run mean & $R^2$ \\
\midrule
1982--2022 & daily     & 10,252 & 4.69   & 111.77 & 2.44  & 0.0007 \\
1982--2022 & monthly   & 713    & 4.97   & 101.19 & 2.55  & 0.0152 \\
1982--2022 & quarterly & 238    & 5.51   & 97.42  & 2.52  & 0.0387 \\
1990--2022 & daily     & 8,256  & 5.56   & 75.23  & 1.77  & 0.0006 \\
\textbf{1990--2022} & \textbf{monthly} & \textbf{574} & \textbf{7.80} & \textbf{61.87} & \textbf{1.76} & \textbf{0.0089} \\
1990--2022 & quarterly & 192    & 6.31   & 74.94  & 1.90  & 0.0273 \\
2000--2022 & daily     & 5,753  & 5.08   & 74.60  & 1.22  & 0.0004 \\
2000--2022 & monthly   & 400    & 7.26   & 62.17  & 1.21  & 0.0062 \\
2000--2022 & quarterly & 134    & 5.14   & 73.78  & 1.23  & 0.0259 \\
2003--2022 & daily     & 5,004  & 364.82 & 72.04  & 86.07 & 0.0000 \\
2003--2022 & monthly   & 348    & 522.79 & 59.78  & 86.01 & 0.0000 \\
2003--2022 & quarterly & 116    & 18.21  & 73.17  & 4.19  & 0.0016 \\
\bottomrule
\end{tabular}
\end{table}

The 2003--2022 rows are the point. Estimated daily over that window the same
estimator returns a mean-reversion half-life of 365 years and a long-run mean of
86\%, with an $R^2$ of zero to four decimal places. Neither figure is a statement
about interest rates. They are what a regression returns when the series it is given
spends most of its length pinned near a bound and exhibits no reversion to identify,
and a model built on them would have been nonsense in a way that no amount of
subsequent calibration would have revealed. The specification used here was chosen
because of this table, and a guard in the configuration rejects any fit implying a
half-life beyond fifty years.

The honest reading extends further than the failed rows, though. The best $R^2$
anywhere in the table is 0.0387, and the chosen specification's is 0.0089. Mean
reversion in the short rate is weakly identified under every specification tried, not
merely the degenerate ones. What makes this tolerable rather than fatal is the result
from Table~\ref{tab:sweep}: the quantity the paper reports is close to insensitive to
this parameter. A weakly identified input that the answer barely depends on is a
different situation from a weakly identified input that drives it. The deposit
assumption is the second kind.

\subsection{The deposit assumption, which dominates everything}

The economic value of equity is the difference between an asset valuation the filings
constrain tightly and a deposit valuation they do not constrain at all. Nothing in
either bank's disclosure states how long its non-maturity deposits are assumed to
remain.

Figure~\ref{fig:heatmap} shows SVB's modelled change in economic value at $+200$
basis points over a grid of 25 deposit betas and 25 uninsured half-lives. Across that
surface the answer runs from $-122.24\%$ to $+21.34\%$. Holding beta at the 0.60 SVB
itself discloses and varying only the half-life, the figure runs from $-120.44\%$ at
six months to $+13.84\%$ at ten years. The same balance sheet, the same securities,
the same curve, and the sign of the exposure is set by an assumption that no filing
contains.

\begin{figure}[htbp]
\centering
\includegraphics[width=0.85\textwidth]{figures/05_deposit_heatmap}
\caption{Modelled change in SVB's economic value of equity at $+200$ basis points,
31~December 2021, across deposit beta and uninsured-deposit half-life. The two
reference points discussed in the text sit inside this surface.}
\label{fig:heatmap}
\end{figure}

Two points on that surface deserve names. Both are defensible, and they disagree by
more than the quantity being measured.

The first is the half-life implied by SVB's own disclosure. Solving for the value that
reproduces the published $-13.9\%$ at $+100$ basis points gives 3.59~years, and at
that value the model returns $-27.61\%$ at $+200$ basis points, against the $-27.70\%$
SVB published. That is the calibration of Section~\ref{sec:comparison}.

The second is a two-year half-life, which is the conservative case used elsewhere in
this paper. At two years the same balance sheet shows $-62.47\%$.

The gap between those two numbers is roughly 35 percentage points of economic value,
which is larger than the entire disclosed exposure at $+200$ basis points. And the
choice between them is not a modelling detail. A 3.59-year half-life on uninsured
deposits is a claim that a depositor base which was 96\% uninsured and concentrated in
venture-backed technology firms would, on average, leave half its balances in place
for three and a half years through a 200 basis point tightening. That is an assumption
about depositor loyalty, and the mechanism in \citet{Drechsler2021} is explicit that
the deposit franchise hedges a long asset book only for as long as the deposits stay.
\citet{Jiang2024} supply the reason the assumption was fragile here specifically:
uninsured depositors are the ones with both the incentive and the capacity to leave.

So the finding is not that SVB's disclosed economic value was wrong. It is that the
disclosed figure was the output of an unstated assumption which, moved to the other
end of a defensible range, more than doubles the measured exposure, and that the
assumption was least likely to hold for precisely this depositor base.

\subsection{Run thresholds}

The run model inherits the same problem in a different form.
Section~\ref{sec:results} reports a threshold of 67.5\% of uninsured deposits for SVB
at the end of 2022, binding through liquidity rather than solvency. Varying secured
borrowing access and the fire-sale discount across 21 combinations moves that
threshold between 22.2\% and 70.5\%. Solvency never binds in the base case, but it
does bind in some cells of that grid, so the correct statement is that liquidity binds
first under the base assumptions and solvency can bind where asset sales are deeply
discounted. Schwab crosses no threshold in either year under any combination tried.

\subsection{The 2024 shock recalibration}

One further variant concerns the measurement convention rather than the balance sheet.
The Basel Committee recalibrated its standardised shocks in 2024, and for USD the
long-rate shock rises from 150 to 225 basis points with effect from 1~January 2026
\citep{BCBS2024}. That post-dates both valuation dates, so it is not the base case,
but it is worth running because the long shock enters the steepener and flattener and
leaves the parallel and short scenarios untouched by construction.

Re-running under the 2024 long shock changes which scenario binds rather than how
badly. For SVB the steepener becomes the worst of the six in all three of its cases,
where the parallel rise held that place before. On the disclosure-implied deposit
assumption the worst-case loss moves from $-27.61\%$ to $-32.59\%$ of base economic
value, and from 36.24\% to 42.78\% of Tier~1 capital. Schwab is unaffected in both
years, because its modelled parallel exposure remains worse than even the strengthened
steepener. No bank-year changes its outlier-test verdict: all five breach the 15\%
threshold under both calibrations.

The implication is narrow and worth stating as such. Under the convention that takes
effect in 2026, a supervisor applying the standardised scenarios to SVB's end-2021
balance sheet would have been directed to a curve rotation rather than a level shift,
and would have measured a loss roughly five percentage points larger. This is a
counterfactual about measurement conventions and not a claim about what SVB's exposure
was.
